\chapter{Software environment and first test}
\label{chap:software}

\section{Prepare the environment}

BrainChip distributes the Akida tools through MetaTF. For the version used in this
manual, the current installation uses \code{metatf==2.19.3} and supports
Python 3.10--3.12 \cite{brainchip_install_2193}.

The repository also includes a reference configuration based on Python 3.11 and
Akida 2.19.1, the version used in the thesis. This second option repeats the smoke
test in a known environment.

Open a terminal in \projectroot, the directory extracted from the ZIP.

On Linux:

\begin{lstlisting}[language=bash]
python3.11 -m venv .venv
source .venv/bin/activate
python -m pip install --upgrade pip
\end{lstlisting}

On Windows, in PowerShell:

\begin{lstlisting}
py -3.11 -m venv .venv
.\.venv\Scripts\Activate.ps1
python -m pip install --upgrade pip
\end{lstlisting}

Create the virtual environment only once. In a new terminal, activate it again
before proceeding.

\section{Install Akida}

Choose one of the following two configurations.

For the current version documented by BrainChip:

\begin{lstlisting}
python -m pip install "metatf==2.19.3"
\end{lstlisting}

To repeat the reference configuration included with the repository:

\begin{lstlisting}
python -m pip install -r requirements/reproducible.txt
\end{lstlisting}

Do not combine both installations in one virtual environment.

Then install the tools included in this repository:

\begin{lstlisting}
python -m pip install -e . --no-deps
python -m pip check
\end{lstlisting}

Check that Akida imports successfully:

\begin{lstlisting}
python -c "import akida; print(akida.__version__); print(akida.devices())"
\end{lstlisting}

Without a connected board, \code{akida.devices()} returning \code{[]} is normal.
With an AKD1000 connected and a working driver, at least one device should appear.

\section{First software test}

Before using the board, run the test model included with the repository:

\begin{lstlisting}
python examples/00_simulator_smoke.py
\end{lstlisting}

The model receives four input values. For

\begin{lstlisting}
[3, 2, 1, 4]
\end{lstlisting}

the expected output is

\begin{lstlisting}
[5, 5]
\end{lstlisting}

The script first runs the model in software, then checks that it can map to a virtual
AKD1000. This last step verifies model compatibility with the board's architecture;
it does not use physical hardware.

You can also run the check provided by the repository's command-line tool:

\begin{lstlisting}
akd1000-bringup verify --id smoke --model models/akida_v1_smoke_test.fbz
\end{lstlisting}

If both tests pass, the Python environment and test model are ready.

\section{Check the host}

Before using the physical board, you can run the repository's diagnostics:

\begin{lstlisting}
akd1000-bringup doctor
\end{lstlisting}

To save the result:

\begin{lstlisting}
mkdir -p evidence
akd1000-bringup doctor --json > evidence/preflight.json
\end{lstlisting}

The tool checks the Akida version, kernel, PCIe connection and detected devices,
among other information. Its output helps locate the fault when the board does not appear.

On a host without an AKD1000, a diagnostic report showing no detected devices is normal.

\section{First execution on the AKD1000}

This test requires PCIe detection of the board and the driver from
\cref{chap:host-driver} to be installed.

First check that Akida detects the device:

\begin{lstlisting}
akida devices
\end{lstlisting}

Then run:

\begin{lstlisting}
python examples/01_hardware_smoke.py
\end{lstlisting}

The script loads the same model used in the software test, maps it to the AKD1000
and runs the known input. The output should again be:

\begin{lstlisting}
[5, 5]
\end{lstlisting}

A \code{PASS} result confirms that the model has loaded, the board is accessible
through Akida and inference has run successfully on the hardware.

\paragraph{If the physical test fails.}
Do not move on to your own model yet. Save the error message and use the procedure
in \cref{chap:diagnostico} to locate the fault.
